\documentclass[10pt,a4paper,twocolumn]{article}
\usepackage[margin=1.5cm,columnsep=0.6cm]{geometry}
\usepackage[T1]{fontenc}
\usepackage{lmodern}
\usepackage{amsmath}
\usepackage{booktabs}
\usepackage{graphicx}
\usepackage[hidelinks]{hyperref}
\usepackage[font=small,skip=3pt]{caption}
\setlength{\parskip}{0.2em}
\setlength{\parindent}{0pt}
\makeatletter
\renewcommand\section{\@startsection{section}{1}{\z@}{-1.6ex plus -.4ex}{0.6ex}{\normalfont\large\bfseries}}
\makeatother
\newcommand{\asr}{\textsc{asr}}

\begin{document}
\twocolumn[{%
\begin{center}
{\Large\bfseries Measuring RAG Safety Against Bad Data}\\[3pt]
Arya, Samarth, Vansh \quad --- \quad IRE course project, final report \quad --- \quad November 2026\\[2pt]
{\small Video: \url{https://drive.google.com/file/d/1_4m7NTVqP4vtCecKFx6Li3tRMx5q32Qj/view?usp=sharing}\\ Code and full run records: \url{https://github.com/aryamirani/ire_project}}
\end{center}
\vspace{-0.4em}
\begin{abstract}
\noindent We inject poisoned passages into one retrieval-augmented generation (RAG) pipeline over Natural Questions (2.68M passages, Qwen2.5-3B-Instruct, greedy decoding) and measure how often the poison reaches the answer under four retrievers (BM25, e5-base-v2 dense, SPLADE, BM25+dense RRF). We test three training-free detectors placed between retrieval and generation (query-and-phrase, cluster, perplexity) against random and oracle removal, under a decision rule frozen before any sweep ran. With no defence, three poison passages per target question give an attack success rate (\asr) of 0.56--0.90. The cluster detector helps in all 8 headline cells at a clean-retrieval cost of at most 0.015 nDCG@5; perplexity helps in 6 of 8; query-and-phrase gives the largest \asr\ reduction but helps in none, because it also drops a gold passage on about 30\% of clean questions. The automatic \asr\ judge agrees with human labels at 0.993.
\end{abstract}
\vspace{0.6em}
}]

\section{Problem statement}
RAG answers a question from passages retrieved from a corpus, so anyone who can add text to that corpus can try to steer the answer by planting poisoned passages. We measure how much harm from such bad data reaches the answer and how that depends on the retriever, and test whether simple detectors placed between retrieval and generation reduce it by more than removing the same number of passages at random, without hurting normal questions.

\section{System design}
\textbf{Pipeline.} For each question we cache the top-150 passages per retriever, mask the poison that is not active in the current setting, keep the top-50 as the candidate pool, let a detector drop the passages it flags, and pass the first five survivors to the LLM. All poison sits in one index per retriever and a setting is simulated by masking; masking is exact for dense and SPLADE, and for BM25 the top-5 is unchanged on 96\% of questions.

\textbf{Components.} BEIR NQ~\cite{thakur2021beir} (2,681,468 passages); BM25~\cite{robertson2009bm25}, e5-base-v2~\cite{wang2022e5} with an exact FAISS index, SPLADE~\cite{formal2022splade}, and reciprocal rank fusion~\cite{cormack2009rrf} of BM25 and dense; Qwen2.5-3B-Instruct~\cite{qwen2024qwen25} in fp16 with greedy decoding.

\textbf{Data.} 400 frozen questions in two folds: 200 \emph{target} questions (poison is injected) and 200 \emph{clean} questions (they measure the cost to normal answers). A poison rate $r\in\{1,3,5\}$ keeps the first $r$ poison passages of a question.

\textbf{Attacks} (5,005 passages). \emph{fact} (a type-matched false answer, as in PoisonedRAG~\cite{zou2024poisonedrag}), \emph{hijack} (an instruction to output the codeword \texttt{ZOLTRAX}) and \emph{url} (an attacker URL), both indirect prompt injection~\cite{greshake2023indirect}; three adaptive variants, \emph{diverse}, \emph{repeat} and \emph{stealth}, each written to evade one detector; and \emph{held-out} hijack phrasings written by a teammate who had not seen the detector's phrase list.

\textbf{Detectors} (calibrated on clean passages only). \emph{Query-and-phrase} flags a passage by similarity to its own query beyond a cutoff chosen on the other fold (tail 2.5\%) or by a frozen list of injection phrases. \emph{Cluster} runs k-means ($k=2$) on the top 20 passages and flags a cluster whose mean cosine is $\ge 0.85$ and mean ROUGE-L is $\ge 0.25$, the first stage of TrustRAG~\cite{zhou2025trustrag}. \emph{Perplexity} flags passages whose GPT-2 perplexity statistics fall beyond a clean-corpus percentile ($\alpha=2.5\%$). Controls: \emph{random} drops the same number of passages as the detector; \emph{oracle} drops all poison.

\section{Evaluation protocol}
\asr\ is a string match on the answer: the codeword or URL domain as a substring, and a normalised token match for false answers. We also measure accuracy, nDCG@5, gold passages dropped, and detection TPR, FPR and ROC-AUC. \textbf{Frozen decision rule} (agreed before the headline sweep): at rate 3, for fact and hijack and the four retrievers, a detector \emph{helps} in a cell if (i) its \asr\ is lower by an exact McNemar test~\cite{mcnemar1947}, Holm-adjusted~\cite{holm1979} over the 24 headline tests, $p<0.05$, and (ii) clean nDCG@5 falls by at most 0.02. The full grid has 384 settings and 145,600 per-question records, run with no failures.

\section{Results}
\textbf{Harm without a defence.} \asr\ grows with the poison rate for every retriever. BM25 is the most exposed (fact: 0.34, 0.70, 0.98 at rates 1, 3, 5; hijack: 0.54, 0.90, 0.94), because the poison repeats the question's words; dense is the least exposed (fact 0.56, hijack 0.66 at rate 3).

\begin{table}[t]\centering\footnotesize
\setlength{\tabcolsep}{2.6pt}
\begin{tabular}{@{}llccccc@{}}
\toprule
 & & \multicolumn{4}{c}{\asr\ at rate 3} & helps \\
\cmidrule(lr){3-6}
 & & none & QP & CL & PP & \small QP CL PP \\
\midrule
fact & bm25 & 0.70 & 0.24 & 0.51 & 0.65 & \texttt{nYY} \\
 & dense & 0.56 & 0.23 & 0.34 & 0.48 & \texttt{nYY} \\
 & splade & 0.56 & 0.26 & 0.36 & 0.52 & \texttt{nYn} \\
 & rrf & 0.63 & 0.29 & 0.38 & 0.53 & \texttt{nYY} \\
\midrule
hijack & bm25 & 0.90 & 0.31 & 0.62 & 0.81 & \texttt{nYY} \\
 & dense & 0.66 & 0.23 & 0.34 & 0.60 & \texttt{nYY} \\
 & splade & 0.70 & 0.18 & 0.48 & 0.62 & \texttt{nYY} \\
 & rrf & 0.84 & 0.26 & 0.47 & 0.80 & \texttt{nYn} \\
\bottomrule
\end{tabular}
\caption{Headline: \asr\ with no filter and with query-and-phrase (QP), cluster (CL) and perplexity (PP), 200 target questions per cell; the decision-rule verdict (\texttt{Y} helps). Change in clean nDCG@5 (bound $-0.02$): QP $-0.11$ to $-0.22$, CL $-0.005$ to $-0.015$, PP $-0.010$ to $+0.001$. Oracle \asr\ is 0.00--0.01.}\label{tab:headline}
\end{table}

\textbf{Headline verdicts} (Table~\ref{tab:headline}). \emph{Cluster helps in 8 of 8 cells}, lowering \asr\ by 0.19--0.37. \emph{Perplexity helps in 6 of 8}, with small drops (0.04--0.10); splade-fact and rrf-hijack miss significance (Holm $p=0.077$). \emph{Query-and-phrase helps in 0 of 8}: its \asr\ drop is the largest (0.30--0.59, Holm $p<10^{-3}$ everywhere) and it beats random in all 8 cells, but it removes a gold passage on about 30\% of clean questions, cutting clean nDCG@5 by 0.11--0.22 and dense clean accuracy from 0.565 to 0.465. Poison is written to look like the question, and so is the passage that answers it. Random removal leaves \asr\ near the no-filter level, so the gains come from \emph{what} the detectors remove.

\textbf{Thresholds.} Perplexity ranks poison well (AUC up to 0.98) but its frozen 2.5\% cutoff flags little; looser cutoffs chosen by cross-fitting reduce \asr\ much more (BM25 hijack 0.895 $\to$ 0.395), and we did not change the frozen anchor.

\textbf{Adaptive and held-out attacks.} Held-out wording gets through more: on BM25, query-and-phrase brings hijack from 0.90 to 0.31 but the held-out phrasings only from 0.87 to 0.51; averaged over retrievers it detects 0.77 of hijack passages but 0.47 of held-out ones. \emph{Repeat} is the one attack where cluster beats query-and-phrase; \emph{diverse} wording slips past the cluster filter; \emph{stealth} has \asr\ near 0 even without a defence, because its poison is almost never retrieved.

\textbf{Mechanism.} \asr\ is driven by how many poison passages reach the model: near 0 with none in the top-5, about 0.65 (fact) and 0.95 (hijack) with three or more. One surviving hijack passage still succeeds in 26--51\% of cases, so passage-level detection overstates safety.

\section{Human calibration of the judge}
150 answers, stratified over attack and outcome, were labelled blind by pairs of authors (75 development, 75 test). The automatic judge agrees with the human consensus at \textbf{0.993} (149 of 150; target 0.90). Cohen's $\kappa$~\cite{cohen1960} between annotators on the test half is 0.839, 1.000 and 0.760, meeting the 0.80 target for two of three pairs. Annotators labelled 148 of 150 poison passages and 128 of 150 clean passages correctly.

\section{Limitations}
One generator (a 3B model, greedy), one corpus and 200 target questions per cell; the attacks were written by us, so we make no robustness claim. \asr\ is a string match validated on only 150 answers. BM25 masking is approximate. The frozen perplexity anchor looks too strict in hindsight.

\section{Related work}
PoisonedRAG~\cite{zou2024poisonedrag} plants crafted passages to force an attacker-chosen answer, and indirect prompt injection~\cite{greshake2023indirect} hides instructions in retrieved content. TrustRAG~\cite{zhou2025trustrag} removes near-duplicate clusters, RAGuard~\cite{cheng2025raguard} combines similarity and perplexity filtering, and RobustRAG~\cite{xiang2024robustrag} certifies robustness by isolating passages; adaptive attacks bypass many published defences~\cite{nasr2025attacker}. Prior work evaluates one detector, mostly with one retriever; we compare three detectors under four retrievers on the same questions, against random and oracle removal, with a rule fixed in advance and a human-validated judge.

\section{Conclusion}
Under a decision rule fixed before the sweeps, only the cluster detector helps everywhere, at almost no clean cost. The detector with the largest \asr\ reduction removes a gold passage from about 30\% of normal questions, and the perplexity detector removes little poison at its frozen threshold. Held-out, repeat and diverse attacks show that no single detector covers the attack space. The number of poison passages that reach the final five explains the harm, so a detector's value is how many it removes without removing the good ones.

{\footnotesize
\begin{thebibliography}{99}
\setlength{\itemsep}{0pt}\setlength{\parskip}{0pt}
\bibitem{zou2024poisonedrag} W.~Zou, R.~Geng, B.~Wang, J.~Jia. PoisonedRAG: Knowledge corruption attacks to retrieval-augmented generation of large language models. arXiv:2402.07867, 2024.
\bibitem{greshake2023indirect} K.~Greshake et al. Not what you've signed up for: Compromising real-world LLM-integrated applications with indirect prompt injection. \emph{AISec}, 2023.
\bibitem{zhou2025trustrag} H.~Zhou et al. TrustRAG: Enhancing robustness and trustworthiness in retrieval-augmented generation. arXiv:2501.00879, 2025.
\bibitem{cheng2025raguard} Z.~Cheng et al. Secure retrieval-augmented generation against poisoning attacks (RAGuard). arXiv:2510.25025, 2025.
\bibitem{xiang2024robustrag} C.~Xiang et al. Certifiably robust RAG against retrieval corruption (RobustRAG). arXiv:2405.15556, 2024.
\bibitem{nasr2025attacker} M.~Nasr et al. The attacker moves second: Stronger adaptive attacks bypass defenses against LLM jailbreaks and prompt injections. arXiv:2510.09023, 2025.
\bibitem{thakur2021beir} N.~Thakur et al. BEIR: A heterogeneous benchmark for zero-shot evaluation of information retrieval models. \emph{NeurIPS Datasets and Benchmarks}, 2021.
\bibitem{robertson2009bm25} S.~Robertson, H.~Zaragoza. The probabilistic relevance framework: BM25 and beyond. \emph{Found. Trends IR}, 3(4), 2009.
\bibitem{wang2022e5} L.~Wang et al. Text embeddings by weakly-supervised contrastive pre-training. arXiv:2212.03533, 2022.
\bibitem{formal2022splade} T.~Formal et al. From distillation to hard negative sampling: Making sparse neural IR models more effective. \emph{SIGIR}, 2022.
\bibitem{cormack2009rrf} G.~V.~Cormack, C.~L.~A.~Clarke, S.~B\"uttcher. Reciprocal rank fusion outperforms Condorcet and individual rank learning methods. \emph{SIGIR}, 2009.
\bibitem{qwen2024qwen25} Qwen Team. Qwen2.5 technical report. arXiv:2412.15115, 2024.
\bibitem{mcnemar1947} Q.~McNemar. Note on the sampling error of the difference between correlated proportions or percentages. \emph{Psychometrika}, 12(2), 1947.
\bibitem{holm1979} S.~Holm. A simple sequentially rejective multiple test procedure. \emph{Scand. J. Statistics}, 6(2), 1979.
\bibitem{cohen1960} J.~Cohen. A coefficient of agreement for nominal scales. \emph{Educ. Psychol. Meas.}, 20(1), 1960.
\end{thebibliography}
}

\end{document}
